\documentclass[11pt]{article}
\usepackage[margin=1in]{geometry}
\usepackage{booktabs}
\usepackage{amsmath,amssymb}
\usepackage{xcolor}
\usepackage{hyperref}
\hypersetup{colorlinks=true, linkcolor=blue!50!black, urlcolor=blue!50!black}
\newcommand{\yes}{\textcolor{green!55!black}{\textbf{verified}}}
\newcommand{\corr}{\textcolor{orange!85!black}{\textbf{corrected}}}

\title{\textbf{Machine-Verification of the HyMeKo Propositions}\\[2pt]
\large Symbolic proofs and property-tests against the real compiler (T-SMC appendix)}
\author{Aiko (Claude Code), for Dr.\ Csaba Hajdu}
\date{2026-06-28 (JST)}

\begin{document}
\maketitle

\begin{abstract}
\noindent The article's four propositions are presented as proof-sketches-by-inspection, with the load-bearing
one (Prop.~1) exiled to supplementary material --- the \#1 weakness a prior review flagged. We attack it
mechanically. The propositions split by mathematical character, so the achievable rigor differs per proposition:
the \emph{algebraic} claims (storage bound, collision bound) admit full \textbf{symbolic proofs} (\textsf{sympy});
the \emph{structural} claims (alias invariance, emitter purity/independence) are \textbf{property-tested against
the real compiler} (the engine's \texttt{parse\_dsl}$\to$\texttt{canonical\_hash} / emitters), which tests the
actual artifact rather than an idealisation. Result: \textbf{three of four propositions are machine-confirmed},
and the fourth surfaces a \textbf{precise, fixable correction} --- Prop.~1's \emph{byte-equal emit} claim fails
under sibling reordering (the emitters carry source-order IDs), although its content-address claim holds on every
one of 800 tested instances. Tooling installed and pinned via \texttt{uv}: \textsf{sympy 1.14, z3 4.16,
hypothesis 6.15, networkx 3.6, pygraphviz 2.0}.
\end{abstract}

\section{Results at a glance}
\begin{table}[h]
\centering
\small
\begin{tabular}{llll}
\toprule
Prop. & Claim & Method & Outcome \\
\midrule
P1 alias invariance & denotation-equal $\Rightarrow$ same IR \emph{and} & Hypothesis vs.\ real compiler, & hash: \yes\ (800/800)\\
 & byte-equal emit & 800 reorderings $\times$ 4 fixtures & emit: \corr\ (0/400)\\
P2 content-address & $h$ pure + collision-resistant & \textsf{sympy} bound + P1 mechanism & \yes\ (+ stated assumption)\\
P3 proj.--emit indep. & emitters pure; $k$ emits, 1 compile & determinism property-test & \yes\ (4/4)\\
P4 storage overhead & $\rho-1=O(\log n/\bar d)\to 1$ & \textsf{sympy} symbolic proof & \yes\ (full proof)\\
\bottomrule
\end{tabular}
\end{table}

\section{P4 --- storage overhead (full symbolic proof)}
From the size equations $|H|_{\mathrm{IR}} = (n+m)c_{\mathrm{rec}} + m\bar d\,c_{\mathrm{inc}}$ and
$|H|_{\mathrm{adj}} = m\bar d\,c_{\mathrm{inc}}$, \textsf{sympy} derives
$\rho = 1 + c\,\frac{n+m}{m\bar d}$ ($c=c_{\mathrm{rec}}/c_{\mathrm{inc}}$), substitutes the assumption
$n\le K m\log n$ to obtain $\rho-1 \le \frac{c(K\log n+1)}{\bar d} = O(\log n/\bar d)$, confirms
$\lim_{\bar d\to\infty}(\rho-1)=0$, confirms strict monotone decrease ($\partial_{\bar d}(\rho-1)<0$), and
reproduces the empirical witness exactly: with $n_{\mathrm{pool}}=m=200,\ c=1$, $\rho-1=2/\bar d$, giving
$\rho=2.00$ at $\bar d=2$ and $\rho=1.01$ at $\bar d=200$. Every step is machine-checked
(\texttt{verification/propositions/p4\_storage\_overhead.py}). This is a proof, not a sketch.

\section{P2 --- content-addressability (symbolic bound + a stated assumption)}
P2 has two parts. (a) \emph{$h$ depends only on the canonical IR} --- the same mechanism as P1, property-tested
below at 800/800. (b) \emph{distinct IRs collide negligibly} --- a birthday bound on the 256-bit digest, proved
symbolically: $P(\text{collision among }N) \le N^2/2^{257}$, so $P\le 2^{-128}$ while $N\le 2^{64.5}$ (work factor
$2^{128}$); numerically $N{=}2^{32}\!\to\!2^{-193}$, $N{=}2^{64}\!\to\!2^{-129}$. \emph{Honest caveats}: the
random-oracle premise for Blake3 is a cryptographic \emph{assumption}, not proved; and a 256-bit digest cannot be
injective (pigeonhole), so equality-by-hash is \emph{practical}, not absolute --- which the article already
concedes.

\section{P1 \& P3 --- against the real compiler}
We compile self-contained fixtures, apply denotation-preserving \emph{sibling reorderings} (the strongest
automatable rewrite), and compare against the engine. Distinct fixtures hash distinctly (fano \texttt{b88715ed},
generic \texttt{6a0c09d8}, k4 \texttt{0a8a903e}, sunflower \texttt{194c71e4}), so the test is non-degenerate.

\paragraph{Verified.} The \textbf{canonical hash is invariant on all 800 reorderings} (4 fixtures $\times$ 200) ---
zero violations: Prop.~1's content-address claim holds against the real compiler. Emitters are
\textbf{deterministic} (4/4): re-compiling the same source in a fresh engine emits byte-identical DOT, so each
emitter is a pure function of its input (Prop.~3's independence then follows --- a pure emitter that ignores the
others can be run $k$ times off one compile).

\paragraph{Corrected.} Prop.~1's \emph{stronger} claim --- denotation-equal sources emit \emph{byte-identical}
text --- \textbf{fails under sibling reordering} (0/400 byte-equal DOT). The cause is precise: the emitters assign
internal identifiers (\texttt{n2}, \texttt{e5}, \dots) in \emph{source-declaration order}, so a reordered source
emits a DOT whose IDs are permuted (not merely the lines --- the ID$\leftrightarrow$label mapping changes). The
graph is isomorphic; the serialisation is not byte-equal. \emph{Only the hash is order-canonicalised; the stored
IR and its emit retain source order.} The article's existing \textsf{prop1\_alias} test uses
\emph{alias-renaming} fixture pairs (where names canonicalise and byte-equality plausibly holds); it does not
exercise sibling reordering, which is why this gap was invisible.

\paragraph{Remedy (for the article).} Three options, in increasing strength: (i) \emph{scope} the byte-equal
claim to alias/import rewrites and state the sibling-reorder guarantee at the canonical-hash level; (ii)
\emph{canonicalise emit IDs} (assign by the canonical sort key before emitting), making byte-equality hold for
reordering too; (iii) \emph{weaken} the claim to ``output equal up to isomorphism / equal canonical hash,'' which
holds universally. Option~(i) is a one-paragraph edit; option~(ii) is the principled fix.

\section{The rigor ladder (kept honest)}
\begin{itemize}\itemsep2pt
\item \textbf{Proved (symbolic):} P4 (storage), P2 collision bound --- \textsf{sympy}, fully rigorous.
\item \textbf{Property-verified (real artifact, falsification over 800 instances):} P1 canonical-hash
  invariance, P3 emitter determinism.
\item \textbf{Assumed (not provable here):} Blake3 random-oracle (P2); sibling reordering taken as
  denotation-preserving \emph{by} the engine's own canonical hash.
\item \textbf{Corrected:} P1 byte-equal emit under sibling reordering.
\item \textbf{z3:} installed, but the four propositions have no natural decidable-SMT formulation that adds over
  \textsf{sympy}+property-testing; its value here is future work (signed-balance gauge properties, bounded
  model-checking the resolver), not these four.
\end{itemize}

\section{What this buys the article}
The appendix moves from ``proof sketches by inspection'' to ``symbolically proved where algebraic, and
property-verified over 800 instances \emph{against the real compiler} where structural'' --- and the exercise
\emph{found a precise, fixable defect} the prose proof had hidden. This directly answers the review's \#1 finding,
and suggests a genuine new angle: a \textbf{machine-checked canonical-IR appendix} (or a short companion artifact)
is a stronger contribution than one more robot fixture. Scripts:
\texttt{verification/propositions/\{p1\_alias\_invariance, p2\_content\_birthday, p3\_emit\_purity,
p4\_storage\_overhead\}.py}.

\end{document}
